\chapter*{专用术语注释表}
\addcontentsline{toc}{chapter}{专用术语注释表}

\section*{符号说明}
\begin{supertabular}{p{3cm}p{10cm}}
$cwnd$ & 拥塞窗口（Congestion Window），TCP发送端在未收到确认的情况下最多可以发送的数据量（字节） \\
$ssthresh$ & 慢启动阈值（Slow Start Threshold），区分慢启动和拥塞避免阶段的临界值（字节） \\
$\alpha$ & 乘性增加因子（Multiplicative Increase Factor），控制拥塞避免阶段窗口增长速度，取值范围$[0.85, 1.30]$，基准值1.10 \\
$\gamma$ & 加性增加因子（Additive Increase Factor），控制窗口的加性增量，默认值为1.0 \\
$\rho$ & 窗口保留因子（Window Retention Factor），拥塞响应后新窗口与原窗口的比值 \\
$\rho_{\text{loss}}$ & 丢包保留因子，默认值0.70 \\
$\rho_{\text{ecn}}$ & ECN保留因子，默认值0.75 \\
$\rho_{\text{timeout}}$ & 超时保留因子，默认值0.50 \\
$\eta_f,\ \eta_b$ & 性能反馈值快线/慢线指数移动平均（EMA）平滑系数，默认值0.15/0.02 \\
$BDP$ & 带宽延迟积（Bandwidth-Delay Product），网络链路能够容纳的最大在途数据量（字节） \\
$MSS$ & 最大报文段长度（Maximum Segment Size），本文取1440字节 \\
$RTT$ & 往返时延（Round-Trip Time），数据包从发送端到接收端再返回的总时间 \\
$W(t)$ & 时刻$t$的拥塞窗口大小 \\
$W_{max}$ & 上次拥塞事件前的最大窗口大小 \\
$K$ & CUBIC算法中恢复到$W_{max}$所需的时间 \\
$C$ & 链路容量或CUBIC缩放因子 \\
$\beta$ & 窗口缩减因子 \\
$\delta$ & 延迟惩罚权重参数 \\
$\lambda$ & 丢包惩罚权重参数 \\
$H$ & Hurst参数，描述流量自相似性 \\
$\bar{Q}$ & 平均队列长度 \\
$\bar{W}$ & 平均等待时间 \\
$\tau$ & 反馈延迟（约等于RTT） \\
\end{supertabular}

\section*{缩略词}
\begin{supertabular}{p{3cm}p{10cm}}
TCP & 传输控制协议（Transmission Control Protocol） \\
ECN & 显式拥塞通知（Explicit Congestion Notification） \\
RTT & 往返时延（Round-Trip Time） \\
BDP & 带宽延迟积（Bandwidth-Delay Product） \\
AIMD & 加性增乘性减（Additive Increase Multiplicative Decrease） \\
MSS & 最大报文段长度（Maximum Segment Size） \\
DCTCP & 数据中心TCP（Data Center TCP） \\
BBR & 瓶颈带宽和往返传播时间（Bottleneck Bandwidth and Round-trip propagation time） \\
CUBIC & 三次函数拥塞控制算法 \\
ACK & 确认（Acknowledgment） \\
SACK & 选择性确认（Selective Acknowledgment） \\
CE & 拥塞经历（Congestion Experienced） \\
ECE & ECN回显（ECN-Echo） \\
CWR & 拥塞窗口缩减（Congestion Window Reduced） \\
RED & 随机早期检测（Random Early Detection） \\
CA & 拥塞算法（Congestion Algorithm） \\
PPO & 近端策略优化（Proximal Policy Optimization） \\
DQN & 深度Q网络（Deep Q-Network） \\
PCC & 面向性能的拥塞控制（Performance-oriented Congestion Control） \\
ZMQ & ZeroMQ消息队列 \\
SDN & 软件定义网络（Software Defined Networking） \\
QoS & 服务质量（Quality of Service） \\
QUIC & 快速UDP互联网连接（Quick UDP Internet Connections） \\
RDMA & 远程直接内存访问（Remote Direct Memory Access） \\
PFC & 优先级流量控制（Priority Flow Control） \\
INT & 带内网络遥测（In-band Network Telemetry） \\
HPCC & 高精度拥塞控制（High Precision Congestion Control） \\
WAN & 广域网（Wide Area Network） \\
Pod & 交换机组（Point of Delivery） \\
EMA & 指数移动平均（Exponential Moving Average） \\
IPC & 进程间通信（Inter-Process Communication） \\
\end{supertabular}
